Le code ci-dessous permet de créer la classe Triangle, entièrement définie par le numéro de ses sommets dans le maillage. Elle comprend également un getter qui renvoie le numéro du sommet demandé. Nous avons également surchargé l'opérateur == afin de facilité les tests.

\begin{codeblock}[triangle.hpp]
#ifndef TRIANGLE_H
#define TRIANGLE_H

class Triangle {

private:
    int n1, n2, n3;

public:
    Triangle(int n1, int n2, int n3);
    Triangle();
    int get(int k);

    // Surcharge de ==
    bool operator==(const Triangle &autre) const;
};


#endif
\end{codeblock}

\begin{codeblock}[triangle.cpp]
#include "triangle.hpp"


Triangle::Triangle(int n1, int n2, int n3){
    this->n1 = n1;
    this->n2 = n2;
    this->n3 = n3;
}


Triangle::Triangle(){
    this->n1 = 0;
    this->n2 = 0;
    this->n3 = 0;
}


int Triangle::get(int k){
    switch (k){
        case 0 : return this->n1;
        case 1 : return this->n2;
        case 2 : return this-> n3;
        default : return -1;
    }
}

// Fonction qui renvoie le min de 3 entiers
int min3(int a, int b, int c) {
    int m = (a < b) ? a : b;
    return (m < c) ? m : c;
}


// Fonction qui renvoie le max de 3 entiers
int max3(int a, int b, int c) {
    int m = (a > b) ? a : b;
    return (m > c) ? m : c;
}


// Surcharge de l'opérateur == pour définir l'égalité de triangles
bool Triangle::operator==(const Triangle &t) const {
    int min1 = min3(n1,n2,n3);
    int max1 = max3(n1,n2,n3);
    int mid1 = n1 + n2 + n3 - min1 - max1;

    int min2 = min3(t.n1,t.n2,t.n3);
    int max2 = max3(t.n1,t.n2,t.n3);
    int mid2 = t.n1 + t.n2 + t.n3 - min2 - max2;

    return min1==min2 && mid1==mid2 && max1==max2;
}
\end{codeblock}